\section{Sharp decision memory for the collision experiment}
\label{sec:v20-physical}
We now apply the continuation spectrum to the same marked experiment
that supplies the microscopic certificate. Both its impossibility
statement and its construction concern the same score, the same private
candidate family, and the same training--validation checkpoint.

\subsection{The marked family and its first two spectral values}
In the all-on row of Section~\ref{sec:v18-microscopic}, a private
width-one candidate has two independent report bits with probabilities
$p,q\in[0,1]$, independently of the fair actual mark. Denote its law by
$b(p,q)$ and put
\[
 a=(1-p)(1-q),\qquad c=pq.
\]
Let $B=\{b(p,q):p,q\in[0,1]\}$. The ideal target $Q_*$ is supported
on $000,001,110,111$, with masses $5/16,3/16,3/16,5/16$, respectively.
On these four atoms the candidate masses are $a/2,a/2,c/2,c/2$.
The elementary continuation constraint implies
\begin{equation}\label{eq:v20-root}
 \sqrt a+\sqrt c\le1,\qquad \min(a,c)\le\tfrac14.
\end{equation}
Indeed the first inequality is Cauchy--Schwarz applied to
$(\sqrt{1-p},\sqrt p)$ and $(\sqrt{1-q},\sqrt q)$; the second follows
from the first. If $a\ge c$, then $c\le1/4$ and
\begin{equation}\label{eq:v20-linear}
 a+3c\le(1-\sqrt c)^2+3c
             =1-2\sqrt c+4c\le1.
\end{equation}
The symmetric inequality holds when $c\ge a$.

\begin{theorem}[A sharp two-state obstruction]
\label{thm:v20-two}
For the ideal marked collision family,
\begin{equation}\label{eq:v20-two}
 \mathfrak V_1(B,Q_*)=\tfrac18,
 \qquad \mathfrak V_2(B,Q_*)=\tfrac38.
\end{equation}
If $\|Q-Q_*\|_{\rm TV}\le\beta$, every reset auditor retaining at most
two states at the training--validation cut has worst-case expected score
at most $3/8+\beta$, regardless of its finite training length and its
upstream or downstream memory. The assertion allows mark-dependent
validation events and stochastic state transitions.
\end{theorem}
\begin{proof}
For one state, use the event $\{000,111\}$. Its score is
$5/8-(a+c)/2\ge1/8$. Conversely, average the two candidates with
$(a,c)=(1,0)$ and $(0,1)$. Their average is uniform on the four supported
atoms and is at total variation distance $1/8$ from $Q_*$. Every fixed
response has average score at most $1/8$, so its minimum over the two
candidates is at most $1/8$. This proves the first equality; the average
is a bound on scores, not an allowed private candidate.

For the upper bound with two states, consider the three candidates
$b_0=b(0,0)$, $b_1=b(1,1)$ and $b_*=b(1/2,1/2)$. A response may be
set to zero outside the support of $Q_*$ without decreasing its payoff
at any candidate. Write its four remaining coordinates as $(x,y,z,t)$.
Its three payoffs are
\begin{align*}
 f_0&=(-3x-5y+3z+5t)/16,\\
 f_1&=(5x+3y-5z-3t)/16,\\
 f_*&=(3x+y+z+3t)/16.
\end{align*}
Since all four coordinates are in $[0,1]$,
\begin{equation}\label{eq:v20-triangle}
 f_0+f_1\le\tfrac14,
 \qquad f_0+f_*\le\tfrac34,
 \qquad f_1+f_*\le\tfrac34.
\end{equation}
Suppose two responses covered all three candidates with payoff strictly
greater than $3/8$. The first inequality forces different responses to
cover $b_0$ and $b_1$. The other two inequalities then force both responses
to have payoff less than $3/8$ at $b_*$. This is a contradiction. Hence
$\mathfrak V_2\le3/8$.

For attainment take the responses, in the same four-atom order,
\[
 h_0=(2/3,0,1,1),\qquad h_1=(1,1,0,2/3),
\]
with zero response elsewhere. Both have target mean $17/24$. Their
payoffs are $17/24-a/3-c$ and $17/24-a-c/3$. If $a\ge c$, the first
is larger and \eqref{eq:v20-linear} makes it at least
$17/24-1/3=3/8$. The other case is symmetric. This proves equality.
For the perturbed target, apply \eqref{eq:v20-spectrumstable} and the
operational upper bound of Theorem~\ref{thm:v20-spectrum}.
\end{proof}

The obstruction has an immediate continuation interpretation. Each
post-training state supplies one validation response. The three witness
candidates in the proof cannot be covered above $3/8$ by two such
responses. No increase of sample size can compensate for losing the
third continuation at that cut. This argument does not use the unrelated
Bernoulli majority calibration.


\begin{lemma}[Feedback completion of the physical lower bound]
\label{lem:v20-feedback}
The upper bound $3/8+\beta$ in Theorem~\ref{thm:v20-two} remains valid
when the auditor may train on any of the original feedback rows and may
select a row--event pair after the decision cut. Here every target row is
the declared feedback pushforward of $Q$, with
$\|Q-Q_*\|_{\rm TV}\le\beta$. In particular, two states cannot evade
the bound by using a different sensor gate during training or validation.
\end{lemma}
\begin{proof}
It suffices to restrict the adversary to an admissible subfamily. For
$p,q\in[0,1]$, let its first emission be $X_1\sim\operatorname{Ber}(p)$,
and, after gate $A$, let its second emission be $AX_2$, with an independent
$X_2\sim\operatorname{Ber}(q)$ and the independent fair actual mark.
These are genuine width-one private rows: the gate is a current input,
and the second emission uses no retained copy of $X_1$.

For any deterministic gate policy $g:\{0,1\}\to\{0,1\}$, the marked
row is the pushforward of $b(p,q)$ under
$T_g(x_1,x_2,w)=(x_1,g(x_1)x_2,w)$. The physical target row is the
pushforward of $Q$ by the same map, since acquisition does not act on the
path. Including the gate in the transcript changes nothing: it is also
a deterministic function of the displayed coordinates. For a selected
row event, the response pulls back to $\1_A\circ T_g$ on the all-on
alphabet. A randomized row--event continuation pulls back to a function
$h_m\in[0,1]^\Omega$. Hence, conditional on cut state $m$, its expected
score on this subfamily is exactly $(Q-b(p,q))h_m$.

Whatever training rows were used, their effect is solely the distribution
of the retained state. The two-state score is thus a convex combination
of two such responses and is bounded above by their envelope. Apply the
three-candidate argument of Theorem~\ref{thm:v20-two}, and then the
uniform target perturbation bound. Since this is a subfamily of the full
private candidate class, it is a valid lower-bound witness for that full
class. Randomized feedback is covered by conditioning on its independent
random table and then averaging the pulled-back response.
\end{proof}

\subsection{A three-response construction}
Define three deterministic events, all ignoring the mark,
\[
 A_D=\{00,11\}\times\{0,1\},\quad
 A_0=\{00\}\times\{0,1\},\quad
 A_1=\{11\}\times\{0,1\}.
\]
Their ideal scores are $1-a-c$, $1/2-a$ and $1/2-c$.

\begin{lemma}[The three-response envelope]
\label{lem:v20-envelope}
For the original private candidate image,
\begin{equation}\label{eq:v20-envelope}
 \min_{p,q}\max\{1-a-c,\tfrac12-a,\tfrac12-c\}=\sqrt2-1.
\end{equation}
In particular $\mathfrak V_3(B,Q_*)\ge\sqrt2-1$.
\end{lemma}
\begin{proof}
By symmetry suppose $a\ge c$. Then $c\le1/4$ and the maximum is
$1-a-c$ for $a\le1/2$, and $1/2-c$ for $a\ge1/2$.
If $a\le1/4$, the first expression is at least $1/2$. If
$1/4\le a\le1/2$, \eqref{eq:v20-root} gives
\[
 a+c\le1-2\sqrt a+2a\le2-\sqrt2,
\]
since the last function is increasing on $[1/4,1/2]$.
If $a\ge1/2$, then
$c\le(1-\sqrt a)^2\le(1-1/\sqrt2)^2=3/2-\sqrt2$.
Both cases give the asserted lower bound. It is attained at
$p=q=1-1/\sqrt2$, where $a=1/2$ and $c=3/2-\sqrt2$.
The symmetric point gives the same value.
\end{proof}

\begin{lemma}[Odd-sample majority selection with a uniform regret bound]
\label{lem:v20-majority}
Let $n=2h-1$ be odd. In $n$ independent candidate trials, let $C_0,C_1$
count report pairs $00,11$. Choose $A_1$ if $C_0\ge h$, choose $A_0$
if $C_1\ge h$, and otherwise choose $A_D$. The first two cases are
mutually exclusive. Its expected ideal score satisfies
\begin{equation}\label{eq:v20-majorityscore}
 g_n(p,q)\ge\sqrt2-1-\frac1{4\sqrt{en}}
                            -\frac14e^{-n/8}.
\end{equation}
Its event label has three possible values at the decision cut.
\end{lemma}
\begin{proof}
Because $a+c\le1$,
\[
 \max\{1-a-c,\tfrac12-a,\tfrac12-c\}
 =1-a-c+(a-1/2)_++(c-1/2)_+.
\]
The selected payoff is
$1-a-c+(a-1/2)\1_{\{C_0\ge h\}}
             +(c-1/2)\1_{\{C_1\ge h\}}$.
Its expected loss from this envelope is exactly the sum, for $u=a,c$,
of $|u-1/2|$ times the probability that the binomial majority selects
the wrong side of $u=1/2$. At $u=1/2$ the contribution is zero.
No independence between $C_0$ and $C_1$ is used.

For $0<u<1/2$ and $k\ge h$, the ratio of the binomial mass at $u$ to
that at $1/2$ is
\[
 [4u(1-u)]^{n/2}\left(\frac{u}{1-u}\right)^{k-n/2}
                  \le[4u(1-u)]^{n/2}.
\]
The fair binomial upper tail is exactly $1/2$ because $n$ is odd.
Reflecting $k$ treats $u>1/2$. The endpoint cases follow directly or by
continuity. Thus the wrong-side probability is at most
\[
 \tfrac12[1-4(u-1/2)^2]^{n/2}
             \le\tfrac12 e^{-2n(u-1/2)^2}.
\]
For $x\ge0$, the maximum of $xe^{-2nx^2}/2$ is
$1/(4\sqrt{en})$, reached at $x=1/(2\sqrt n)$.
At least one of $a,c$ is at most $1/4$ by \eqref{eq:v20-root}; that
coordinate has $|u-1/2|\in[1/4,1/2]$ and contributes at most
$e^{-n/8}/4$. Bound the other coordinate by the unrestricted maximum
and use Lemma~\ref{lem:v20-envelope}.
\end{proof}

\subsection{Simultaneous exact widths and the full acquisition schedule}
The preceding construction admits more compression than storing its two
raw counts. Write $0,1,*$ for the three report-pair categories $00,11$
and the two off-diagonal pairs. For $n=2h-1$, let $F_n$ be the terminal
event-label function of this word. At a trial cut $t$, let $r=n-t$ and
let $(u,v)$ be the two prefix counts. Associate the residual key
\begin{equation}\label{eq:v20-key}
 \mathcal R_t(u,v)=
 \begin{cases}
  \mathsf L,&u\ge h,\\
  \mathsf R,&v\ge h,\\
  (\min(h-u,r+1),\min(h-v,r+1)),&u,v<h.
 \end{cases}
\end{equation}
The symbols $\mathsf L,\mathsf R$ mean the constant decisions $A_1,A_0$.
In a pair key, a coordinate $r+1$ means that this majority is no longer
reachable. The pair $(r+1,r+1)$ represents the constant decision $A_D$.

\begin{theorem}[Exact profile of the physical audit rule]
\label{thm:v20-profile}
The keys in \eqref{eq:v20-key} are exactly the distinct continuation
residuals of $F_n$. Their number at trial cut $t$ is
\begin{equation}\label{eq:v20-profile}
 r_t(F_n)=
 \begin{cases}
  \binom{t+2}{2},&0\le t\le h-1,\\[2pt]
  \binom{n-t+3}{2},&h\le t\le n.
 \end{cases}
\end{equation}
This is the minimum simultaneous width profile for exact computation of
this rule, for deterministic and stochastic machines alike. Its maximum
is $\binom{h+2}{2}$ at $t=h$.

For $n=399$, the maximum trial-cut width is $20301$, and the terminal
training cut has exactly three states. Acquiring each report pair one
bit at a time, discarding its training mark, and then executing independent
candidate and target validations gives a deterministic machine with at
most $40602<2^{16}$ states at every scheduled epoch. One additional
ignored training trial may be inserted to give the same $400$ training
slots and $402$ total trials as the earlier certificate. It changes
neither the payoff nor these width bounds.
\end{theorem}
\begin{proof}
If a majority has already been reached, its label cannot change. Otherwise,
the first coordinate of the pair key is exactly the number of additional
$0$ symbols needed, with the unreachable threshold clipped to $r+1$;
the second coordinate has the analogous meaning. These data determine
the suffix function. Two different reachable pair keys can be distinguished
by a suffix containing only one of the two symbols and enough $*$ symbols
to fill the remaining positions: use a count between the two different
thresholds. The other majority cannot be reached by this suffix. The
constant-majority keys are distinct from every pair key because an all-$*$
suffix returns $A_D$ for the latter. Thus the keys are exactly the
residuals.

For $t<h$, no majority has been reached and every threshold is at most
$r+1$. The distinct pairs $(u,v)$ with $u+v\le t$ give
$\binom{t+2}{2}$ keys. For $t\ge h$, the two constant-majority keys
are reachable. The remaining keys are precisely
\[
 (i,j)\in\{1,\ldots,r+1\}^2,\qquad i+j\ge r+1.
\]
To see surjectivity, use $u=h-i$, $v=h-j$; these counts are nonnegative,
less than $h$, and sum to at most $t$ because $2h-t=r+1$.
To see necessity before clipping, use the same inequality on $h-u,h-v$;
if either coordinate is clipped, the inequality remains automatic.
There are
$(r+1)^2-(r-1)r/2$ such pairs, with the expression valid also for
$r=0,1$. Adding the two constant keys gives
$(r+2)(r+3)/2$. This proves \eqref{eq:v20-profile}. The finite-valued
realization theorem, Theorem~\ref{thm:v20-weighted}, proves both the
simultaneous realization and stochastic minimality for the specified rule.
The two polynomial pieces of the profile give its maximum as stated.

For an explicit transition definition, order the reachable keys
lexicographically, putting $\mathsf L,\mathsf R$ first. For each key
choose the lexicographically first prefix-count pair representing it.
After the next category, increment the appropriate count and apply
$\mathcal R_{t+1}$. Equality of residual functions implies equality of
the shifted residuals, so this definition is independent of the chosen
representative. The representative is part of the time-indexed transition
table, not another stored register. The clock supplies $t$ but no sample
history. All these maps are deterministic and rational.

At a bit-level epoch, the state consists of a residual key and, only
between the two report bits, the first bit. It therefore has at most
$2\max_t r_t$ values. On reading the second bit, determine the category,
update the key and discard both bits. The mark is read and discarded;
all-on controls and resets are prescribed by the free schedule. After
training, keep just one of the three event labels. During validation,
keep that label, the completed candidate-event indicator, and a first-bit
buffer when needed. At most $3\cdot2\cdot2=12$ states suffice there;
a completed target pair emits its event indicator minus the stored
candidate indicator. The final score has three possible values. These
phase-specific alphabets may be embedded in a single alphabet of size
$40602$ because the schedule clock identifies the phase.

The same trial-cut residuals remain distinct for the complete signed-score
function, not merely for an externally emitted event name. Any two of
$A_0,A_1,A_D$ have different indicator functions. Take an off-diagonal
candidate-validation pair, which belongs to none, and a target-validation
pair distinguishing the two events. The resulting scores differ. Thus
passing to the validation suffix does not merge the residuals already
counted.
\end{proof}

\begin{theorem}[Matched physical decision complexity and stable certification]
\label{thm:v20-physical}
For the two-sphere collision family of
Theorem~\ref{thm:v18-microscopic}, let $d\in[9/10,11/10]$ and
$0<\sigma\le1/24$, and let $Q_{d,\sigma}$ be its all-on marked target.
At the charged training--validation cut of the reset interface,
\begin{equation}\label{eq:v20-kappa}
 \kappa_{2/5}(B,Q_{d,\sigma})=3.
\end{equation}
More concretely, no auditor retaining at most two states at this cut can
have expected score greater than $2/5$ uniformly against all private
width-one candidates, at any finite training length, even with feedback-row selection as in
Lemma~\ref{lem:v20-feedback}. One rational auditor
with three states at this cut does so with $399$ informative training
trials and two validation trials. With a redundant ignored training trial,
it uses exactly $402$ total trials, at most $40602$ states at every
scheduled epoch, and the exact trial-cut profile
\eqref{eq:v20-profile} on its informative training segment.

For fixed positive $\sigma$ and $d_k\to d$, the same machine changes its
expected score by at most
\begin{equation}\label{eq:v20-transportscore}
 \epsilon_k=\min\left\{1,
 \frac{\sqrt{|d_k-d|+2700|d_k-d|^2}}{2\sigma}\right\}.
\end{equation}
The private, hidden and visible deficiencies, the no-collision control,
and the joint microscopic, likelihood, posterior, innovation and global
bounded entropic backward convergence retain all conclusions of
Theorems~\ref{thm:v18-microscopic} and~\ref{thm:v19-global}.
\end{theorem}
\begin{proof}
Theorem~\ref{thm:v18-microscopic} discharges the physical hypotheses and
proves $\|Q_{d,\sigma}-Q_*\|_{\rm TV}<1/300$, uniformly on the stated
parameters. Theorem~\ref{thm:v20-two} therefore gives the universal
cut-width-two upper bound
\[
 \mathfrak V_2(B,Q_{d,\sigma})<\frac38+\frac1{300}<\frac25.
\]
Lemma~\ref{lem:v20-feedback} makes this a bound for the full original
feedback-row interface, not only all-on training and not only the
constructed rule. The upper construction uses all-on trials and works
for every private candidate because only its two all-on parameters enter.

Use the deterministic selector of Lemma~\ref{lem:v20-majority} with
$n=399$, but validate against the actual physical target. For each
selected event, changing only the target from $Q_*$ loses less than
$1/300$. Its physical expected score is consequently greater than
\begin{equation}\label{eq:v20-certifiedscore}
 L_{399}:=\sqrt2-1-\frac1{4\sqrt{399e}}
                       -\frac14e^{-399/8}-\frac1{300}>\frac25.
\end{equation}
Here is a rational verification of the last strict inequality, independent
of decimal evaluation. The bounds $e>2$, $\sqrt{798}>28$,
$\sqrt2>707/500$, and $399/8>49$ imply
\[
 L_{399}>\frac{207}{500}-\frac1{112}-\frac1{300}-2^{-51}
                >\frac25,
\]
since $207/500-1/112-1/300-2/5=73/42000>2^{-51}$.
The number $L_{399}$ is approximately $0.4032891$; the rational inequalities,
not this approximation, establish the certificate. Its state realization
was proved in Theorem~\ref{thm:v20-profile}. The general operational
classification now yields \eqref{eq:v20-kappa} with both inequalities
proved for this same physical task.

Changing $d$ leaves the blind candidate source and all its training laws
unchanged. The moving-collision estimate and the common-preparation
Gaussian theorem bound the target variation by $\epsilon_k$.
For any fixed response its target mean changes by at most this amount;
averaging over the unchanged training rule proves
\eqref{eq:v20-transportscore}. There is no multiplier equal to the number
of training trials. The same physical estimates verify the hypotheses of
Theorem~\ref{thm:v19-global}; the observation-level density and coupling
versions used there are given in Section~\ref{sec:v20-versions}.
This transports the same certificate and the bounded conditional-law
consumer along the changing microscopic paths. The no-collision regime
$d\le2/5$ has the vanishing deficiencies already proved in
Theorem~\ref{thm:v18-microscopic}; it is not included in the positive-gap
parameter interval.
\end{proof}

\subsection{What the two exact complexity statements measure}
Equation~\eqref{eq:v20-kappa} minimizes over every auditor at one
specified physical decision cut. Equation~\eqref{eq:v20-profile}
minimizes all trial-cut widths simultaneously for exact implementation
of the stated $402$-trial rule. These are different optimizations.
The first is a task-level optimum, not an artifact of that rule. The
second explains the whole online realization of the same rule, not a
claim that every successful auditor must have $20301$ states somewhere.
The sixteen-bit bound includes bit acquisition and validation but does
not bound the description length of time-dependent transition tables.
It is an implementation upper bound for peak memory. Independent
persistent randomization is charged; fresh randomness from the reset
experiments is not a hidden cross-cut memory channel.

The former two-count $400$-sample construction and its nineteen-bit bound
remain in Section~\ref{sec:v19-certificate}. They offer a different
approximation to the full marked discrepancy. The new three-response
rule does not reproduce that rule or attain its oracle envelope. It
instead solves the specified physical certification task with an exactly
minimal decision-cut alphabet and an explicit full continuation profile.
The score is an expectation, not a confidence guarantee from a single
realized validation pair.
